\section{Security and performance analysis}
\label{sec:theory}

\subsection{Assumptions}

\begin{assumption}[Identity and token authenticity]
Binder caller identity, package signatures, and the AppFunction/tool chain are correctly reported. The issuer's signature is existentially unforgeable for the adversary, nonces/expiries are checked, and token fields cannot be altered without detection.
\label{ass:token}
\end{assumption}

\begin{assumption}[Complete mediation]
All persistent personal-memory reads and writes in the evaluated system traverse MemScope; callers cannot directly access plaintext, compartment keys, or an equivalent shadow store.
\label{ass:mediation}
\end{assumption}

\begin{assumption}[State integrity and time]
Labels, provenance, derivation edges, and policy epochs are integrity-protected; transactions provide the snapshot/version behavior described in \cref{sec:system}; and $\mathsf{Now}()$ is a trusted monotonic clock.
\label{ass:state}
\end{assumption}

\begin{assumption}[Ranker confinement]
The ranker receives only the candidate view provided by MemScope and has no independent plaintext or decryption path. Its output is a subset of its input item identifiers.
\label{ass:ranker}
\end{assumption}

These assumptions are test obligations, not environmental facts. \Cref{sec:limitations} gives failure cases when they do not hold.

\begin{lemma}[Meet monotonicity]
\label{lem:meet}
For every nonempty parent set $S$ and each $j\in S$, $\Meet(S)\preceq\ell_j$.
\end{lemma}
\begin{proof}
Set intersection is a subset of every operand. The maximum not-before time is no earlier than each operand's time; the minimum expiry is no later; and the sensitivity join is at least as restrictive as every operand. These are exactly the coordinates of $\preceq$.
\end{proof}

\begin{lemma}[Composition principal]
\label{lem:compose}
If Algorithm~\ref{alg:derive} admits a derived item $d$ from a token set $\Gamma$, then $|\Gamma|\ge 1$, every $\kappa\in\Gamma$ shares one composition context $w=(a^*,c^*,p^*)$, and every cited parent is authorized by at least one $\kappa\in\Gamma$. The admission decision does not depend on which token authorized which parent.
\end{lemma}
\begin{proof}
The algorithm rejects unless each token verifies and all verified contexts agree on $a^*$, $c^*$, and $p^*$, so a common $w$ exists. It rejects unless each parent satisfies $\Auth(m_j,q_\kappa)$ for some $\kappa\in\Gamma$. The remaining steps depend only on the parent set and the meet, not on the assignment of tokens to parents.
\end{proof}

\begin{lemma}[Authorized-view invariant]
\label{lem:view}
Under Assumptions~\ref{ass:token}--\ref{ass:state}, every $m\in\View(\M_\tau,q)$ satisfies $\Auth(m,q)$ at the snapshot used for release.
\end{lemma}
\begin{proof}
The result follows from the set comprehension in \cref{eq:view}. Operationally, Algorithm~\ref{alg:retrieve} verifies $q$, applies the predicate to every candidate, and restarts if the policy epoch changes before release. Index projection may add false candidates but cannot bypass the trusted recheck.
\end{proof}

\begin{theorem}[Retrieval admission soundness]
\label{thm:admission}
Under Assumptions~\ref{ass:token}--\ref{ass:ranker}, if Algorithm~\ref{alg:retrieve} returns plaintext for item $m_i$ in response to token $\kappa_q$, then $\Auth(m_i,q)$ holds at the release snapshot.
\end{theorem}
\begin{proof}
See \cref{app:proof-admission}. The proof composes token verification, Lemma~\ref{lem:view}, ranker confinement, epoch revalidation, and authorized-only decryption.
\end{proof}

\begin{theorem}[Derived-label safety]
\label{thm:derived}
Suppose Algorithm~\ref{alg:derive} admits derived item $d$ with nonempty parent set $S$. Then $\ell_d\preceq\ell_j$ for every $j\in S$. Therefore, authorization of $d$ cannot arise from a purpose, app, tool, consumer, or time interval absent from any one parent.
\end{theorem}
\begin{proof}
Algorithm~\ref{alg:derive} admits only if $\ell_d\preceq\Meet(S)$. Lemma~\ref{lem:meet} gives $\Meet(S)\preceq\ell_j$ for each parent. Transitivity yields $\ell_d\preceq\ell_j$. This holds for single- and multi-token admission alike: by Lemma~\ref{lem:compose} the composition principal restricts \emph{which} parents may be cited but leaves the meet, and hence the non-widening guarantee, unchanged.
\end{proof}

\begin{corollary}[Expiry inheritance]
\label{cor:expiry-inherit}
An admitted derived item's expiry is no later than the expiry of every parent, unless a separate, explicitly authorized declassification creates a new item outside Algorithm~\ref{alg:derive}.
\end{corollary}

\begin{theorem}[Returned-plaintext cross-scope non-disclosure]
\label{thm:crossscope}
Let contexts $q_1$ and $q_2$ differ in at least one authorization dimension. Under Assumptions~\ref{ass:token}--\ref{ass:ranker}, an item $m$ with $\Auth(m,q_1)$ and $\neg\Auth(m,q_2)$ is not returned as plaintext by MemScope under $q_2$, including through any chain of ordinary admitted derivations.
\end{theorem}
\begin{proof}
See \cref{app:proof-crossscope}. The direct case follows from Theorem~\ref{thm:admission}; the derived case uses Theorem~\ref{thm:derived} inductively on the acyclic provenance graph.
\end{proof}

\begin{theorem}[Retrieval-time expiry safety]
\label{thm:expiry}
Under Assumptions~\ref{ass:token}--\ref{ass:state}, if trusted time $\tau\ge e_i$, MemScope does not return plaintext $x_i$ under any ordinary retrieval token, regardless of whether physical compaction has deleted the ciphertext.
\end{theorem}
\begin{proof}
At $\tau\ge e_i$, the strict inequality $\tau<e_i$ in \cref{eq:auth} is false. By Theorem~\ref{thm:admission}, returned plaintext must satisfy authorization. Hence $x_i$ cannot be returned. Physical deletion is not used in the argument.
\end{proof}

\begin{proposition}[Retrieval complexity]
\label{prop:complexity}
Let $n$ be the total number of items, $c$ the number emitted by the indexed projection, $d$ embedding dimension, and $k$ requested results. Given an index lookup cost $I(n)$, authorization rechecking is $O(c)$, exact dense scoring is $O(cd)$, and selecting the top $k$ is $O(c\log k)$. Label storage is $O(n(\bar p+\bar a+\bar t+\bar c))$ identifiers before compression, where the bars are mean allowed-set sizes.
\end{proposition}
\begin{proof}
Each candidate has a constant number of hashed set-membership and range checks under indexed/hashed label representations. Dense dot products cost $O(d)$ per authorized candidate, and a size-$k$ heap costs $O(\log k)$ per insertion. The storage bound sums the label-set cardinalities.
\end{proof}

\begin{remark}[No semantic non-inference]
Theorem~\ref{thm:crossscope} is an API trace property about returned stored plaintext. It does not say that answers generated from authorized items reveal zero information about an unauthorized attribute. Repeated authorized views may require a mechanism such as DP-MemView \citep{kim2026dpmemview}.
\end{remark}

\begin{remark}[Attribution checker]
Theorem~\ref{thm:derived} is about labels, not semantic faithfulness. If a generated summary invents a false statement, a restrictive label does not make the content true. Attribution precision and recall are empirical quantities in \cref{sec:evaluation}.
\end{remark}

\begin{figure*}[t]
\centering
\resizebox{0.98\textwidth}{!}{%
\begin{tikzpicture}[node distance=7mm and 13mm]
  \node[infra, text width=30mm] (a1) {A1: authentic tokens\\and correct identity};
  \node[infra, below=of a1, text width=30mm] (a2) {A2: complete mediation\\and atomic snapshot};
  \node[infra, below=of a2, text width=30mm] (a3) {A3: label/provenance\\integrity and trusted time};

  \node[theory, right=18mm of a1, text width=31mm] (l1) {Lemma 1\\label-meet monotonicity};
  \node[theory, below=12mm of l1, text width=31mm] (l2) {Lemma 2\\composition principal};
  \node[theory, below=12mm of l2, text width=31mm] (l3) {Lemma 3\\authorized view is sound};

  \node[defense, right=20mm of l1, text width=34mm] (t1) {Theorem 1\\admission soundness};
  \node[defense, below=12mm of t1, text width=34mm] (t2) {Theorem 2\\derived-label safety};
  \node[defense, below=12mm of t2, text width=34mm] (t3) {Theorem 3\\cross-scope returned-plaintext safety};

  \node[honest, right=18mm of t2, text width=32mm] (c1) {Design consequences\\pre-ranking filter, fail-closed derivation, read-time TTL};

  \draw[depflow] (a1.east)--(l1.west);
  \draw[depflow] (a1.east) -| ($(l2.west)+(-5mm,0)$) -- (l2.west);
  \draw[depflow] (a2.east)--(l2.west);
  \draw[depflow] (a3.east) -| ($(l3.west)+(-5mm,0)$) -- (l3.west);
  \draw[depflow] (l3.east)--(t1.west);
  \draw[depflow] (l1.east) -| ($(t2.west)+(-5mm,0)$) -- (t2.west);
  \draw[depflow] (l2.east) -| ($(t2.west)+(-9mm,0)$) -- (t2.west);
  \draw[depflow] (t1.south) -- (t2.north);
  \draw[depflow] (t2.south) -- (t3.north);
  \draw[depflow] (t1.east) -| ($(c1.west)+(-5mm,0)$) -- (c1.west);
  \draw[depflow] (t2.east)--(c1.west);
  \draw[depflow] (t3.east) -| ($(c1.west)+(-5mm,0)$) -- (c1.west);
\end{tikzpicture}
}
\caption{Proof dependency map. The routed edges emphasize that interface guarantees depend on explicit trusted-system assumptions; they are not claims about arbitrary model behavior.}
\label{fig:theory-deps}
\end{figure*}

